
%%%%%%%%%%%%%%%%%%%%%%%%%%%%%%%%%%%%%%%%%%%%%%%%%%%%%%%%%%%%%
\subsection{模型的缺点}
\begin{itemize}[itemindent=2em]
\item \textbf{振打峰值估算粗糙：}半经验公式中系数0.1为经验值，缺乏标定数据；分钟级数据无法直接观测秒级峰值，估算不确定性较大。
\item \textbf{电场间耦合未建模：}假设$H_1$认为各级独立，但前电场振打释放的粉尘会重返气流进入后电场，振打瞬间级间独立性不成立。
\item \textbf{$C_{out}$外推不确定性：}从历史$C_{out}\approx 50$外推到$10$甚至$5\;\text{mg/Nm}^3$跨两个数量级，Deutsch参数拟合虽有物理约束，但数据欠定下参数不确定性仍较大。贝叶斯95\%置信区间覆盖率仅3\%，因限幅致残差非正态，须谨慎解读。
\item \textbf{振打过频系数$r$不可辨识：}$C_{out}$限幅使$r$无法从数据辨识，当前取$r=0.5$为工程先验。$r$敏感性分析表明最优电耗对$r$不敏感（$r=0.3\sim1.0$下$P^*$几乎相同），但$r$对最优振打周期位置有影响。
\item \textbf{后电场振打参数不可辨识：}贝叶斯估计显示$\alpha_3$、$\alpha_4$的95\% CI宽度达46--47\%且贴下界，后电场振打衰减参数辨识困难。
\end{itemize}